\section{Hướng tiếp cận dự kiến}

\subsection{Các vấn đề của bài toán sau khi khảo sát các công trình liên quan}

\subsection{Hướng tiếp cận của đề tài}

Để giải quyết các vấn đề trên, đề tài lựa chọn mô hình YOLOv8n làm bộ phát hiện
phương tiện. YOLOv8n là phiên bản có kích thước nhỏ của họ YOLOv8, cho phép đạt
tốc độ suy luận cao trong khi vẫn duy trì độ chính xác phù hợp cho các ứng dụng
giao thông thời gian thực. Mô hình được fine-tune trên bộ dữ liệu UA-DETRAC nhằm
nâng cao khả năng nhận diện các loại phương tiện xuất hiện trong môi trường giao
thông thực tế.

Bên cạnh đó, đề tài sử dụng thuật toán ByteTrack để thực hiện theo dõi đa đối
tượng. Khác với nhiều phương pháp theo dõi truyền thống chỉ sử dụng các phát
hiện có độ tin cậy cao, ByteTrack tận dụng cả các phát hiện có độ tin cậy thấp
trong quá trình liên kết đối tượng giữa các khung hình. Nhờ đó, thuật toán giúp
giảm hiện tượng mất dấu, hạn chế thay đổi định danh đối tượng và nâng cao độ
chính xác của quá trình đếm phương tiện.

Hệ thống đề xuất bao gồm ba giai đoạn chính: phát hiện phương tiện bằng
YOLOv8n, theo dõi phương tiện bằng ByteTrack và đếm phương tiện dựa trên quỹ
đạo di chuyển của đối tượng qua vùng hoặc đường đếm được xác định trước. Cách
tiếp cận này hướng tới mục tiêu cân bằng giữa độ chính xác và tốc độ xử lý,
đáp ứng yêu cầu giám sát giao thông tự động và cung cấp dữ liệu gần thời gian
thực cho các ứng dụng trong ITS.

\subsection{Mô hình YOLOv8}

YOLOv8 được Ultralytics phát hành vào ngày 10 tháng 1 năm 2023, mang đến những
cải tiến đáng kể về độ chính xác và tốc độ xử lý so với các phiên bản trước
đó. Kế thừa những ưu điểm của các thế hệ YOLO trước, YOLOv8 được bổ sung nhiều
đặc điểm và tối ưu hóa mới trong kiến trúc mạng, giúp nâng cao khả năng trích
xuất đặc trưng và hiệu quả phát hiện đối tượng. Nhờ những cải tiến này, YOLOv8
trở thành một lựa chọn phù hợp cho nhiều bài toán phát hiện đối tượng trong
các ứng dụng thị giác máy tính khác nhau.

Về mặt kiến trúc, YOLOv8 được xây dựng dựa trên ba thành phần chính: Backbone,
Neck và Detection Head. Thành phần Backbone sử dụng kiến trúc CSPDarknet cải
tiến đảm nhiệm vai trò trích xuất đặc trưng đa tỉ lệ từ ảnh đầu vào thông
qua các lớp tích chập được tối ưu hóa nhằm học các đặc trưng phân cấp bao gồm
đặc trưng mức thấp như kết cấu, cạnh và đặc trưng ngữ nghĩa cấp cao. Các đặc
trưng này sau đó được chuyển đến thành phần Neck để thực hiện hợp nhất thông
tin đa tỉ lệ. Thành phần Neck sử dụng kiến trúc PANet cải tiến giúp tăng cường
khả năng trao đổi thông tin giữa các mức đặc trưng khác nhau, từ đó nâng cao
hiệu quả phát hiện đối tượng có kích thước khác nhau. Thành phần Head chịu
trách nhiệm sinh ra các dự đoán cuối cùng, bao gồm tọa độ hộp bao, độ tin cậy,
và nhãn lớp. Một cải tiến đáng chú ý của YOLOv8 là việc áp dụng cơ chế
anchor-free thay cho cơ chế anchor-based được sử dụng trong các phiên bản
YOLO trước. Cách tiếp cận này giúp đơn giản quá trình dự đoán, giảm số lượng
tham số và cải thiện khả năng thích ứng của mô hình đối với đối tượng có kích
thước khác nhau \cite{yaseen2024yolov8}.

Ngoài những cải tiến về kiến trúc, YOLOv8 còn được tối ưu hóa nhằm đạt sự cân
bằng giữa độ chính xác và tốc độ xử lý. Ultralytics cung cấp nhiều biến thể
khác nhau bao gồm YOLOv8n, YOLOv8s, YOLOv8m, YOLOv8l và YOLOv8x, cho phép
người dùng lựa chọn mô hình phù hợp với yêu cầu về hiệu năng và tài nguyên
phần cứng.

\subsubsection{Biến thể YOLOv8n}

Trong nghiên cứu này, biến thể YOLOv8n được lựa chọn làm mô hình phát hiện đối
tượng chính. YOLOv8n là phiên bản có quy mô nhỏ nhất trong họ YOLOv8, được
thiết kế nhằm giảm yêu cầu về tài nguyên tính toán trong khi vẫn duy trì hiệu
năng phát hiện đối tượng ở mức cao. Với số lượng tham số ít hơn so với các
biến thể còn lại, YOLOv8n phù hợp cho các hệ thống nhúng, thiết bị biên và các
ứng dụng yêu cầu xử lý thời gian thực.

YOLOv8n đạt được kích thước mô hình nhỏ gọn thông qua việc sử dụng các lớp
tích chập được tối ưu hóa và giảm số lượng tham số trong mạng. Kích thước mô
hình chỉ khoảng 2 MB ở định dạng INT8 và khoảng 3.8 MB ở định dạng FP32, giúp
giảm yêu cầu về bộ nhớ lưu trữ và bộ nhớ truy cập trong quá trình suy luận.
Đặc điểm này khiến YOLOv8n trở thành lựa chọn phù hợp cho các thiết bị nhúng,
hệ thống IoT và các nền tảng di động có giới hạn về năng lượng và tài nguyên
phần cứng \cite{yaseen2024yolov8}.

Ngoài ra, YOLOv8n hỗ trợ triển khai trên nhiều nền tảng tăng tốc suy luận như
ONNX Runtime và TensorRT, cho phép tận dụng hiệu quả tài nguyên phần cứng và
giảm độ trễ khi xử lý thời gian thực \cite{yaseen2024yolov8}.

Nhờ sự kết hợp giữa kiến trúc gọn nhẹ, tốc độ cao và khả năng triển khai linh
hoạt, YOLOv8n được lựa chọn làm mô hình nền tảng cho hệ thống phát hiện đối
tượng được đề xuất trong nghiên cứu này.

\input{bytetrack}
